\section{3차시 --- LED 2개로 표정 만들기}
\label{sec:ch03}

본 차시는 LED 2개를 이용하여 알비의 "눈빛" 즉 감정 표현을 만든다.
동시 점등, 교대 점등(윙크), 빠른 깜빡임(놀람), 그리고 학습자 스스로
설계하는 감정 표현으로 이어진다. 두 개의 출력 장치를 \emph{동시에}
다루는 경험을 통해, 코드가 하드웨어 여러 개를 병렬적으로 조작할 수
있다는 사고로 확장된다.

\subsection{미션 1 --- 알비의 눈이 반짝!}
\label{subsec:ch03-m1}
\missionmeta{LED 2개}{WINK}

\begin{storybox}
\sayares{일어나, 알비! 반짝반짝 깜빡이게 만들어볼까?}
\sayalbi{삐릿- 안녕하세요! 양쪽 눈이 동시에 깜빡이면 살아있는 것처럼 보여요!}
\end{storybox}

\subsubsection*{학습 목표}
\begin{itemize}[leftmargin=1.4em]
  \item LED 1번, 2번 동시에 켜기
  \item 1초 대기 후 동시에 끄기
  \item 무한 반복으로 계속 깜빡이게 구현
\end{itemize}

\begin{labbox}{샘플 코드 --- 양쪽 눈 동시 깜빡임}
\begin{minted}{python}
# 양쪽 눈 동시 깜빡이기
while True:
    led_on(1, 1.0)  # 왼쪽 눈
    led_on(2, 1.0)  # 오른쪽 눈
    time.sleep(1)
    led_off(1)
    led_off(2)
    time.sleep(1)
\end{minted}
\end{labbox}

\subsection{미션 2 --- 알비가 윙크해요!}
\label{subsec:ch03-m2}
\missionmeta{LED 2개}{WINK}

\begin{storybox}
\sayalbi{삐릿! 아레스, 저 윙크할 수 있어요! 한쪽 눈씩 번갈아 깜빡이는 거예요!}
\sayares{왼쪽, 오른쪽 눈이 번갈아 켜지면 윙크처럼 보이겠는데? 해보자!}
\end{storybox}

\subsubsection*{학습 목표}
\begin{itemize}[leftmargin=1.4em]
  \item 왼쪽 눈(LED1) 켜기 → 1초 → 끄기
  \item 오른쪽 눈(LED2) 켜기 → 1초 → 끄기
  \item 무한 반복으로 알비가 윙크하듯 구현
\end{itemize}

\begin{labbox}{샘플 코드 --- 윙크 코딩}
\begin{minted}{python}
# 윙크 코딩!
while True:
    led_on(1, 1.0)   # 왼쪽 윙크
    time.sleep(1)
    led_off(1)
    led_on(2, 1.0)   # 오른쪽 윙크
    time.sleep(1)
    led_off(2)
\end{minted}
\end{labbox}

\subsection{미션 3 --- 알비가 깜짝 놀랐어요!}
\label{subsec:ch03-m3}
\missionmeta{LED 2개}{WINK}

\begin{storybox}
\sayares{알비! 갑자기 운석이 날아오면 어떻게 돼?}
\sayalbi{삐릿!! 깜짝 놀라죠! 놀라면 두 눈이 막 빠르게 깜빡이잖아요!}
\end{storybox}

\subsubsection*{학습 목표}
\begin{itemize}[leftmargin=1.4em]
  \item 두 눈 동시에 0.1초 간격 빠른 깜빡이기
  \item 10번 빠르게 깜빡인 후 잠깐 멈추기
  \item 반복문으로 놀람 표현 완성
\end{itemize}

\begin{labbox}{샘플 코드 --- 운석 경보 빠른 깜빡임}
\begin{minted}{python}
# 운석 경보! 빠른 깜빡임
while True:
    for i in range(10):
        led_on(1, 1.0)
        led_on(2, 1.0)
        time.sleep(0.1)
        led_off(1)
        led_off(2)
        time.sleep(0.1)
    time.sleep(0.5)  # 잠깐 멈춤
\end{minted}
\end{labbox}

\subsection{미션 4 --- 나만의 알비 눈빛 만들기}
\label{subsec:ch03-m4}
\missionmeta{LED 2개}{WINK}

\begin{storybox}
\sayares{알비, 이제 직접 원하는 눈빛을 만들어봐!}
\sayalbi{삐릿♪ 기쁠 때, 슬플 때, 생각할 때... 각각 다른 눈빛으로 감정을 표현해줘요!}
\end{storybox}

\subsubsection*{학습 목표}
\begin{itemize}[leftmargin=1.4em]
  \item 기쁨: 두 눈 빠르게 반짝 (0.2초)
  \item 슬픔: 두 눈 천천히 깜빡 (1.5초)
  \item 생각 중: 한쪽 눈만 깜빡
  \item 나만의 감정 눈빛 하나 추가 창작
\end{itemize}

\begin{labbox}{샘플 코드 --- 기쁨과 슬픔 표현}
\begin{minted}{python}
# 기쁨 표현
for i in range(5):
    led_on(1, 1.0); led_on(2, 1.0)
    time.sleep(0.2)
    led_off(1); led_off(2)
    time.sleep(0.2)

# 슬픔 표현
for i in range(3):
    led_on(1, 1.0); led_on(2, 1.0)
    time.sleep(1.5)
    led_off(1); led_off(2)
    time.sleep(1.5)
\end{minted}
\end{labbox}

\begin{notebox}{창작 활동}
미션 4는 학습자 자유 창작 미션이다. 4가지 감정 외에도 학습자가 직접
"화남"(빨갛게 빠르게), "졸림"(점점 느려짐) 등 자신만의 표현을
설계하도록 격려한다.
\end{notebox}
